% COMS 6998, Fall 2026 -- P3: Evaluation Plan & Initial Results
% Based on the course proposal, P1, and P2 templates and the official ISCA 2026 IEEE class.
% Overleaf: upload the ZIP, select main.tex, and use pdfLaTeX.
% Local: latexmk -pdf main.tex
%
% Replace the guidance and bracketed fields with your project content.
% Write two pages, excluding references; retain the supplied formatting.
% Optional bibliography instructions appear at the end of this file.

\documentclass[conference,10pt,letterpaper]{IEEEtran}
\usepackage[T1]{fontenc}
\usepackage{cite}
\usepackage{graphicx}
\usepackage{booktabs}
\usepackage{tabularx}
\usepackage{flushend}
\usepackage[hyphens]{url}
\usepackage[hidelinks]{hyperref}

\newcommand{\field}[1]{\textit{[#1]}}

\title{\field{Project Title}}
\author{\IEEEauthorblockN{\field{Student 1 (UNI)}, \field{Student 2 (UNI)}, \field{Student 3 (UNI)}}
\IEEEauthorblockA{COMS 6998: AI-Native Computing $\mid$ Fall 2026\\
P3: Evaluation Plan \& Initial Results $\mid$ Due November 25, 2026, 11:59 PM ET}}

\begin{document}
\maketitle
\thispagestyle{plain}
\pagestyle{plain}

\noindent\textbf{Format.} Write \textbf{two pages}, excluding references.
Figures and tables count toward the two pages.
Use the four sections below and replace the guidance and bracketed fields
with your project content.

\smallskip
\noindent\textbf{Submission.} Submit this report as a PDF.

\section{Project Goal and Updates}
Briefly state your project goal, the research question, and the main
claim you intend to test. Summarize any updates to your scope or direction
since Milestone~2, including changes made in response to midterm and P2
feedback, and explain the reasons for these changes.

\section{Evaluation Setup}
Describe the evaluation setup you will use for your main results and
briefly justify your choices. Depending on your project, this may include:
\begin{itemize}
\item \textbf{Workloads:} the workloads, datasets, models, or design tasks
you evaluate on, including at least one held-out workload, configuration,
or system scale.
\item \textbf{Baselines:} at least two meaningful baselines, with a
classical or non-agentic baseline where applicable, and what each
represents.
\item \textbf{Metrics:} task quality or success, reported together with
the relevant systems metrics.
\item \textbf{Platform:} the hardware, simulators, models, and tools you
use.
\end{itemize}
Organize this section in the way that best fits your project, for example
as paragraphs, a list, or a table.

\section{Initial Results}
Present initial comparisons between your method and the baselines on the
setup above. Include at least one figure or table, stating the workload,
baselines, metrics, configuration, and number of runs. Explain what the
results suggest, whether they support your intended claim so far, and any
results that differ from your expectations.

\section{Planned Ablations and Next Steps}
List the ablations or controlled interventions you plan to run. For each,
state the component you will remove or vary, the part of your claim it
tests, and the outcome you expect. Identify a failure or boundary case
you plan to study. Outline your plan for \textbf{Milestone~4 on Dec~7}:
main results, at least one ablation, a failure or boundary case, and
repeated measurements or uncertainty where relevant.

% Optional references: add your sources to references.bib, cite them in the
% text using \cite{your_key}, and uncomment the two lines below.
% \bibliographystyle{IEEEtran}
% \bibliography{references}

\end{document}
